
\phantomsection
\section*{Глава 1. Анализ предметной области}
\addcontentsline{toc}{section}{Глава 1. Анализ предметной области}
\stepcounter{section}

\subsection{Обзор существующих программных продуктов по теме работы}

Предметная область проекта связана с автоматизацией работы частной медицинской клиники. К основным процессам такой организации относятся регистрация пациентов, запись на приём, ведение расписания врачей, оформление результатов посещения, хранение медицинских данных и предоставление пациенту доступа к информации о проведённом лечении.

В традиционном варианте часть этих процессов выполняется вручную или с использованием разрозненных программных средств. Администратор принимает звонки и фиксирует запись, врач ведёт медицинские записи в отдельной системе или документе, пациент узнаёт информацию о визитах через администратора. Такая схема имеет несколько недостатков: дублирование данных, риск ошибки при выборе времени приёма, сложность поиска медицинской информации, отсутствие единой истории посещений и повышенная нагрузка на сотрудников.

Разработка веб-приложения позволяет объединить основные процессы в единой информационной системе. Пациент получает возможность самостоятельно записываться на приём и просматривать сведения о визитах через личный кабинет, врач - работать с медицинскими картами и результатами приёмов, а администратор - управлять пользователями, расписанием, услугами и записями. В результате повышается прозрачность работы клиники, сокращается количество ручных операций и улучшается качество обслуживания пациентов.

Для определения функциональных требований были рассмотрены существующие программные продукты, предназначенные для автоматизации медицинских организаций и цифрового взаимодействия с пациентами. В качестве аналогов выбраны Инфоклиника.RU, Medesk, ONDOC и Cliniccards. Эти решения относятся к одной предметной области, но отличаются назначением: одни ориентированы на комплексное управление клиникой, другие - на личный кабинет пациента или на клиники определённого профиля. Поэтому сравнение выполнено в формате «критерий - платформа», что позволяет увидеть, какие функции необходимо включить в разрабатываемое приложение.


{\scriptsize
\sloppy
\setlength{\LTcapwidth}{\textwidth}
\setlength{\tabcolsep}{2pt}
\renewcommand{\arraystretch}{1.12}
\begin{longtable}{|>{\raggedright\arraybackslash}p{0.18\textwidth}|>{\raggedright\arraybackslash}p{0.175\textwidth}|>{\raggedright\arraybackslash}p{0.175\textwidth}|>{\raggedright\arraybackslash}p{0.175\textwidth}|>{\raggedright\arraybackslash}p{0.175\textwidth}|}
\caption{Анализ существующих платформ по функциональным критериям}\label{tab:platform_analysis}\\
\hline
\textbf{Критерий} & \textbf{Инфоклиника.RU} & \textbf{Medesk} & \textbf{ONDOC} & \textbf{Cliniccards} \\ \hline
\endfirsthead
\tabcont{\ref{tab:platform_analysis}}
\textbf{Критерий} & \textbf{Инфоклиника.RU} & \textbf{Medesk} & \textbf{ONDOC} & \textbf{Cliniccards} \\ \hline
\endhead
Назначение платформы & Медицинская информационная система и электронная регистратура для клиники~\cite{infoclinica}. & Облачная медицинская информационная система и CRM для управления клиникой~\cite{medesk}. & Пациентский сервис и личный кабинет, ориентированный на хранение медицинских данных~\cite{ondoc}. & CRM-система для стоматологических клиник и управления пациентской базой~\cite{cliniccards}. \\ \hline
Онлайн-запись на приём & Поддерживается через электронную регистратуру: пациент выбирает специалиста и время. & Поддерживается как часть расписания, CRM и карточки пациента. & Поддерживается со стороны пациента и зависит от подключения клиники к сервису. & Используется вместе с расписанием, но сценарии ориентированы преимущественно на стоматологию. \\ \hline
Личный кабинет пациента & Позволяет просматривать сведения о посещениях, результатах и рекомендациях. & Реализован частично: основной акцент сделан на внутренних процессах клиники. & Является ключевой функцией: пациент получает доступ к медицинской информации и взаимодействию с клиникой. & Имеются элементы работы с карточкой пациента, но личный кабинет пациента не является основной функцией продукта. \\ \hline
Электронная медицинская карта & Поддерживается как часть комплексной медицинской системы. & Поддерживает протоколы осмотров, историю записей и результаты исследований. & Даёт пациенту доступ к накопленным медицинским сведениям и документам. & Поддерживается с учётом стоматологической специфики, включая историю лечения пациента. \\ \hline
Управление расписанием врачей & Расписание связано с электронной регистратурой и используется для записи на приём. & Расписание применяется для планирования загрузки врачей и управления визитами. & Возможности зависят от интеграции с информационной системой клиники. & Расписание используется для планирования приёмов и ведения пациентской базы. \\ \hline
Ролевая модель доступа & Предусмотрены разные уровни доступа для сотрудников клиники, врача и пациента. & Предусмотрены роли сотрудников и механизмы разграничения процессов. & Основной пользователь - пациент; внутренняя ролевая модель зависит от интеграции с клиникой. & Предусмотрены роли сотрудников клиники и доступ к рабочим разделам CRM. \\ \hline
Работа с медицинскими документами & Возможен просмотр результатов анализов, заключений и рекомендаций. & Поддерживаются лабораторные результаты, медицинские документы и данные приёмов. & Поддерживается хранение медицинских данных и документов пациента. & Возможна работа с данными пациента и файлами в рамках CRM-системы. \\ \hline
Администрирование клиники & Включает управление пациентами, расписанием, приёмами и медицинскими процессами. & Включает CRM-функции, управление клиникой, пациентской базой и коммуникациями. & Ограничено, так как продукт ориентирован прежде всего на пациентский сценарий. & Включает управление пациентами, расписанием и рабочими процессами стоматологической клиники. \\ \hline
Специализация решения & Подходит для клиник разного профиля, где нужна электронная регистратура и МИС. & Подходит для частных клиник, которым требуется единая система управления процессами. & Подходит для клиник, которые хотят предоставить пациентам цифровой доступ к медицинской информации. & Подходит прежде всего для стоматологических клиник и профильных медицинских центров. \\ \hline
Открытость архитектуры для анализа & Низкая: коммерческая закрытая система. & Низкая: готовый облачный коммерческий сервис. & Низкая: внутренняя архитектура и логика интеграций не раскрываются. & Низкая: закрытая CRM-система с готовыми сценариями работы. \\ \hline
\end{longtable}
}

Анализ существующих платформ показывает, что цифровые сервисы для клиник обычно объединяют несколько обязательных направлений: онлайн-запись, управление расписанием, ведение электронной медицинской карты, работу с медицинскими документами и разграничение прав доступа. Инфоклиника.RU и Medesk ближе всего к комплексным медицинским информационным системам, ONDOC делает акцент на пациентском личном кабинете, а Cliniccards демонстрирует специализированный подход к ведению карточек пациентов и расписания в стоматологической клинике.

По результатам сравнения можно выделить функции, которые целесообразно учитывать при проектировании веб-приложения для частной клиники. Система должна поддерживать самостоятельную запись пациента на приём, хранение истории посещений, работу врача с медицинской картой, администрирование расписания и защиту данных с помощью ролевой модели. При этом рассмотренные аналоги являются закрытыми коммерческими продуктами, поэтому для курсового проекта важно не копировать готовое решение, а самостоятельно спроектировать структуру данных, пользовательские сценарии и архитектуру приложения.

По итогам анализа к разрабатываемому приложению предъявляются следующие требования:

\begin{enumerate}
\item Наличие нескольких ролей пользователей: пациент, врач и администратор;
\item Регистрация и авторизация пользователей;
\item Управление расписанием врачей и доступными временными интервалами;
\item Создание, изменение и отмена записи на приём;
\item Ведение медицинской карты пациента и истории посещений;
\item Загрузка, просмотр и выгрузка файлов, связанных с медицинскими документами;
\item Личный кабинет пациента с просмотром записей и медицинской информации;
\item Административный интерфейс для управления пользователями, врачами, услугами и записями.
\end{enumerate}

\subsection{Анализ программных инструментов разработки веб-приложений}

При выборе программных инструментов необходимо учитывать специфику предметной области. Приложение для частной клиники работает с персональными и медицинскими данными, поэтому важны надёжная аутентификация, разграничение прав доступа, корректная работа с базой данных, защита форм, удобная разработка CRUD-интерфейсов и возможность дальнейшего развёртывания на сервере.

Для выбора технологической основы рассмотрены несколько распространённых стеков веб-разработки. Сравнение выполнено не по отдельным языкам программирования, а по целостным наборам инструментов, так как серверная часть, шаблонизация, база данных и средства интерфейса должны совместно обеспечивать реализацию требований проекта.


{\scriptsize
\sloppy
\setlength{\LTcapwidth}{\textwidth}
\setlength{\tabcolsep}{2pt}
\renewcommand{\arraystretch}{1.12}
\begin{longtable}{|>{\raggedright\arraybackslash}p{0.145\textwidth}|>{\raggedright\arraybackslash}p{0.205\textwidth}|>{\raggedright\arraybackslash}p{0.205\textwidth}|>{\raggedright\arraybackslash}p{0.185\textwidth}|>{\raggedright\arraybackslash}p{0.14\textwidth}|}
\caption{Анализ существующих стеков разработки веб-приложений}\label{tab:stack_analysis}\\
\hline
\textbf{Стек} & \textbf{Состав и назначение} & \textbf{Преимущества для проекта клиники} & \textbf{Ограничения и риски} & \textbf{Итоговая оценка} \\ \hline
\endfirsthead
\tabcont{\ref{tab:stack_analysis}}
\textbf{Стек} & \textbf{Состав и назначение} & \textbf{Преимущества для проекта клиники} & \textbf{Ограничения и риски} & \textbf{Итоговая оценка} \\ \hline
\endhead
Python + Flask + PostgreSQL + Bootstrap & Flask отвечает за маршрутизацию, обработку запросов, шаблоны Jinja2 и подключение расширений. PostgreSQL хранит пользователей, врачей, расписание, записи, медицинские карты и документы. Bootstrap используется для форм, таблиц и адаптивной вёрстки~\cite{python_docs,flask_docs,postgres_docs,bootstrap_docs}. & Стек даёт понятную архитектуру и позволяет явно реализовать роли, маршруты, модели данных, CRUD-интерфейсы и загрузку файлов. Flask не навязывает лишнюю структуру, поэтому хорошо подходит для учебного проектирования и поэтапной реализации. PostgreSQL удобен для связанных медицинских данных. & Часть механизмов нужно подключать и настраивать самостоятельно: авторизацию, миграции, защиту форм, проверку прав доступа и обработку файлов. Требуется дисциплина в организации структуры проекта. & Оптимальный вариант для проекта на Python: простой, гибкий и достаточный по функциональности. \\ \hline
Python + Django + PostgreSQL & Django включает MVC/MVT-структуру, ORM, административную панель, систему пользователей, формы, миграции и встроенные механизмы безопасности~\cite{django_docs}. & Быстро создаются модели, админ-панель и типовые CRUD-операции. Стек хорошо подходит для крупных информационных систем с большим количеством сущностей и строгой структурой. & Для курсового проекта может быть избыточным: значительная часть логики скрыта во встроенных механизмах. Сложнее показать самостоятельное проектирование маршрутов и пользовательских сценариев. & Надёжный промышленный вариант, но менее наглядный для демонстрации собственной архитектуры. \\ \hline
Node.js + Express + React/Vue + PostgreSQL или MongoDB & Express реализует серверное API, React или Vue формируют клиентское SPA-приложение, база данных хранит доменные сущности. Обычно стек предполагает разделение frontend и backend~\cite{node_docs,express_docs,react_docs}. & Подходит для интерактивного интерфейса, динамического личного кабинета, API и последующего масштабирования клиентской части. Удобен, когда требуется сложный frontend. & Увеличивает объём работы: нужно разрабатывать API, клиентское состояние, сборку frontend, авторизацию токенами и обмен данными между частями приложения. Для базового прототипа это усложняет реализацию. & Подходит для SPA, но избыточен для отчётного прототипа с серверными шаблонами. \\ \hline
PHP + Laravel + MySQL/PostgreSQL & Laravel предоставляет MVC-структуру, маршруты, шаблоны Blade, ORM Eloquent, миграции, middleware, валидацию форм и средства аутентификации~\cite{laravel_docs}. & Удобен для классических веб-приложений с личными кабинетами, формами, административной частью и CRUD-интерфейсами. Имеет развитую экосистему готовых инструментов. & Требует использования PHP и знания экосистемы Laravel. Если проект ориентирован на Python, переход на Laravel не даёт принципиального преимущества и усложняет сопровождение. & Практичная альтернатива, но не соответствует выбранному языку реализации. \\ \hline
Java + Spring Boot + PostgreSQL & Spring Boot используется для серверной части, REST-контроллеров, сервисного слоя, безопасности, работы с БД через JPA/Hibernate и построения корпоративных приложений~\cite{spring_docs}. & Подходит для сложных систем, где важны строгая типизация, многоуровневая архитектура, масштабируемость, тестируемость и развитые механизмы безопасности. & Имеет высокий порог входа, требует больше конфигурации и более сложной структуры проекта. Для учебного прототипа частной клиники может быть слишком тяжёлым. & Надёжный, но избыточный стек для заданного объёма работы. \\ \hline
ASP.NET Core + PostgreSQL/SQL Server & ASP.NET Core позволяет создавать MVC-приложения и API, подключать Identity, Entity Framework Core и работать с реляционными базами данных~\cite{aspnet_docs}. & Хорошо подходит для корпоративных веб-систем, поддерживает строгую архитектуру, производительность и развитые инструменты безопасности. & Требует использования C\# и отдельной экосистемы разработки. Для проекта, ориентированного на Python, усложняет реализацию и не совпадает с планируемым стеком. & Сильный промышленный вариант, но нецелесообразен для выбранной реализации. \\ \hline
\end{longtable}
}

По результатам анализа выбран стек Python + Flask + PostgreSQL + Bootstrap. Он соответствует задачам курсового проекта и позволяет реализовать приложение без избыточного усложнения архитектуры. Flask обеспечивает маршрутизацию, шаблоны и возможность подключать только необходимые расширения: авторизацию пользователей, ORM для работы с базой данных, обработку форм, миграции и загрузку файлов. PostgreSQL подходит для реляционной модели предметной области, где пользователь, пациент, врач, расписание, запись на приём, посещение и медицинский документ связаны между собой. Bootstrap и базовые технологии HTML, CSS и JavaScript позволяют подготовить адаптивный интерфейс личного кабинета, расписания, медицинской карты и административных страниц.

Выбор Flask обоснован тем, что он даёт возможность явно показать структуру веб-приложения: маршруты, модели, формы, шаблоны, роли пользователей и логику доступа. В отличие от Django, он не скрывает значительную часть архитектуры за готовыми механизмами; в отличие от SPA-стеков на Node.js, не требует отдельной разработки сложного frontend-приложения; по сравнению с Laravel, Spring Boot и ASP.NET Core лучше соответствует выбранному языку и учебному объёму проекта. Следовательно, стек Python + Flask является наиболее рациональным вариантом для проектирования и последующей реализации веб-приложения частной клиники.

В качестве инструмента для обеспечения кроссплатформенности и упрощения развертывания веб-приложения обоснован выбор технологии контейнеризации Docker. Использование Docker позволяет упаковать приложение клиники, его зависимости (Flask-расширения, СУБД) и начальную конфигурацию в единый изолированный контейнер, который гарантирует идентичную работу системы на любом компьютере или хост-сервере независимо от установленной операционной системы и локальных библиотек. Это минимизирует риски несовместимости версий интерпретатора Python, исключает конфликты пакетов, ускоряет развертывание и существенно упрощает масштабирование информационной системы медицинского учреждения в процессе промышленной эксплуатации.

\subsection{Формулировка цели и задач работы}

Целью курсового проекта является разработка веб-приложения для частной клиники, обеспечивающего онлайн-запись пациентов на приём, ведение медицинских карт, работу с расписанием врачей и предоставление пользователям личных кабинетов с разграничением прав доступа.

Для достижения поставленной цели в рамках пояснительной записки необходимо решить следующие задачи:

\begin{enumerate}
\item Провести анализ предметной области и рассмотреть существующие программные продукты для автоматизации клиник;
\item Определить функциональные требования к веб-приложению на основе анализа аналогов и задания на курсовой проект;
\item Выполнить полный анализ существующих стеков веб-разработки и обосновать выбор Python + Flask;
\item Проанализировать целевую аудиторию и роли пользователей;
\item Описать функциональность приложения через диаграмму вариантов использования и пользовательские истории;
\item Разработать ER-диаграмму, логическую и физическую схемы базы данных;
\item Подготовить макеты страниц личных кабинетов и основных пользовательских сценариев.
\end{enumerate}

Итогом выполнения курсового проекта должно стать описание и проектная основа веб-приложения, демонстрирующего основные процессы цифрового взаимодействия частной клиники с пациентами. Подготовленные результаты позволяют перейти к программной реализации системы на стеке Python + Flask, закрепить навыки проектирования веб-систем, работы с базами данных, моделирования пользовательских ролей и оформления пояснительной записки в соответствии с требованиями ГОСТ~7.32-2017~\cite{gost_report}.
